\documentclass{article}  % Define la clase del documento.

% Paquetes de idioma y codificación
\usepackage[utf8]{inputenc}
\usepackage[T1]{fontenc}
\usepackage[spanish]{babel}  % Ajusta el idioma del documento a español.

% Paquete de geometría para configurar márgenes y tamaño de papel
\usepackage[letterpaper, margin=3cm]{geometry}

% Paquetes de tipografía
\usepackage{mathptmx}    % Usa Times New Roman como fuente.
\usepackage{microtype}   % Mejora la justificación del texto.

% Paquetes para manejo de colores y gráficos
\usepackage{xcolor}      % Define y utiliza colores.
\usepackage{graphicx}    % Permite la inserción de imágenes.
\usepackage{tikz}        % Creación de gráficos vectoriales.

% Configuración de enlaces y referencias cruzadas
\usepackage{hyperref}
\hypersetup{
    colorlinks   = true,
    linkcolor    = darkblue,
    citecolor    = black,
    filecolor    = blue,
    urlcolor     = blue
}

% Paquetes para la mejora visual de tablas y figuras
\usepackage{booktabs}    % Para tablas de alta calidad.
\usepackage{float}       % Controla la posición de figuras y tablas.

% Paquete para la personalización de códigos fuente
\usepackage{listings}
\lstset{
    literate=
    {á}{{\'a}}1 {é}{{\'e}}1 {í}{{\'i}}1 {ó}{{\'o}}1 {ú}{{\'u}}1
    {Á}{{\'A}}1 {É}{{\'E}}1 {Í}{{\'I}}1 {Ó}{{\'O}}1 {Ú}{{\'U}}1
    {ñ}{{\~n}}1 {Ñ}{{\~N}}1 {ü}{{\"u}}1 {Ü}{{\"U}}1,
    backgroundcolor=\color{backcolour},
    commentstyle=\color{codegreen},
    keywordstyle=\color{codepurple},
    numberstyle=\tiny\color{codegray},
    stringstyle=\color{red},
    basicstyle=\ttfamily\small,
    breakatwhitespace=false,
    breaklines=true,
    captionpos=b,
    keepspaces=true,
    numbers=left,
    numbersep=5pt,
    showspaces=false,
    showstringspaces=false,
    showtabs=false,
    tabsize=2,
    language=TeX,
    morecomment=[l]\#,
    frame=single,
    rulecolor=\color{black}
}

% Definición de colores al estilo Visual Studio Code
\definecolor{darkblue}{rgb}{0.0, 0.0, 0.55}  % Enlaces
\definecolor{codegreen}{rgb}{0.25, 0.49, 0.48}  % Comentarios
\definecolor{codegray}{rgb}{0.5, 0.5, 0.5}  % Números y anotaciones
\definecolor{codepurple}{rgb}{0.58, 0, 0.82}  % Palabras clave
\definecolor{backcolour}{rgb}{0.95, 0.95, 0.92}  % Fondo de código

% Configuraciones de párrafo y matemáticas
\usepackage{amsmath}
\usepackage{parskip}    % Espaciado entre párrafos.
\usepackage{ragged2e}   % Justificación mejorada.

% Configuración de secciones y encabezados
\usepackage{titlesec}
\titleclass{\part}{top} % Make part like a class
\titleformat{\part}[display]
  {\normalfont\huge\bfseries\centering}{\thepart}{20pt}{\Huge}
\titlespacing*{\part}{172.5pt}{-60pt}{10pt}
\titleformat{\part}
  {\normalfont\huge\bfseries}{}{0pt}{}

% Asegúrate de usar esto para mantener el estilo en las páginas de las partes
\titleformat{\part}[display]
  {\normalfont\huge\bfseries}{}{0pt}{}
  [\thispagestyle{fancy}] % Aplica el estilo fancy a las páginas de las partes

% Configuración de encabezados y pies de página personalizados
\usepackage{fancyhdr}
\pagestyle{fancy}
\fancyhf{}
\fancyhead[L]{\raisebox{0.20cm}{\textbf{Tarea 1 - Métodos y Técnicas de Construcción}}}
\fancyhead[R]{\raisebox{0.1cm}{\includegraphics[width=0.25\linewidth]{LOGO_UNIVERSIDAD.jpg}}}
\fancyhead[C]{\rule{\textwidth}{0.6pt}}
\fancyfoot[C]{\rule{\textwidth}{0.6pt}}
\fancyfoot[R]{\raisebox{-1.5\baselineskip}{\thepage}}
\renewcommand{\headrulewidth}{0pt}
\renewcommand{\footrulewidth}{0pt}

% Configuración avanzada de geometría
\geometry{
  top=3.5cm, % Aumenta el espacio en la parte superior para subir el encabezado
  bottom=2.5cm,
  headheight=2.5cm % Aumenta la altura del encabezado si es necesario
}

% Configuracion de bibliografia
\usepackage{natbib}
\bibliographystyle{unsrtnat}  % Puedes cambiarlo por `unsrtnat`, `abbrvnat`, etc.

\begin{document}
%----------------------------------------------------------------------------------------
% PORTADA
%----------------------------------------------------------------------------------------
\begin{titlepage}%Inicio de la carátula, solo modificar los datos necesarios
\newcommand{\HRule}{\rule{\linewidth}{0.5mm}} 
\center 
%----------------------------------------------------------------------------------------
%	ENCABEZADO
%----------------------------------------------------------------------------------------
\includegraphics[width=10cm]{LOGO_UNIVERSIDAD.jpg}\\ % Si esta plantilla se copio correctamente, va a llevar la imagen del logo de la facultad.OBS: Es necesario incluir el paquete: graphicx
\vspace{3cm}
%----------------------------------------------------------------------------------------
%	SECCION DEL TITULO
%----------------------------------------------------------------------------------------
\HRule \\[0.4cm]
{ \huge \bfseries Tarea 1}\\[0.4cm] % Titulo del documento
{ \huge \bfseries Métodos y Técnicas de Construcción}\\[0.4cm] % Titulo del documento
\HRule \\[1.5cm]
 \vspace{5cm}
%----------------------------------------------------------------------------------------
%	SECCION DEL AUTOR
%----------------------------------------------------------------------------------------
\begin{flushright}
    { 
    \textbf{Alumnos:} \\
    Felipe Vicencio\\
    Lukas Wolff Casanova\\
    \vspace{0.5cm}
    \textbf{Profesor:}\\
    José Tramon Cardenas\\
}
\end{flushright}
\vspace{1cm}
%----------------------------------------------------------------------------------------
%	SECCION DE LA FECHA
%----------------------------------------------------------------------------------------
{\large \textbf{\today}}\\[2cm] % El comando \today coloca la fecha del dia, y esto se actualiza con cada compilacion, en caso de querer tener una fecha estatica, reemplazar el \today por la fecha deseada
\end{titlepage}
%----------------------------------------------------------------------------------------
%  INDICE
%----------------------------------------------------------------------------------------
\newpage
\thispagestyle{empty} % Deshabilita el número de página en la página del índice
\tableofcontents
\thispagestyle{plain} % Deshabilita el encabezado en la página del índice
\thispagestyle{empty} % Deshabilita el número de página en la página del índice
\newpage

%\newpage
%\thispagestyle{empty}
%\listoffigures 
%\thispagestyle{plain} % Deshabilita el encabezado en la página del índice %
%\thispagestyle{empty}
%\newpage
%----------------------------------------------------------------------------------------
%ACÁ EMPIEZA EL INFORME
\setcounter{page}{1}
%----------------------------------------------------------------------------------------
\section{Introducción}

La tarea consiste en realizar la cubicación y programación de algunas obras involucradas en la construcción del Aeropuerto y evaluar los plazos para cada actividad indicada. Específicamente, se trata de un proyecto para un paso bajo nivel en una nueva vía de circulación denominada Diego Barros Ortiz. Las partidas incluyen actividades relacionadas con excavaciones, muros de contención, pavimentación y la construcción de cajones de hormigón armado. Esta tarea no solo implica un cálculo preciso de volúmenes, dimensiones y tiempos, sino que también tiene un impacto real en la planificación y ejecución de obras civiles, dado que una adecuada cubicación y evaluación de plazos son esenciales para cumplir con los requisitos técnicos, legales y financieros de proyectos de infraestructura. 
\\\\
El proceso de cubicación y planificación es crucial para determinar la viabilidad y el éxito de una obra civil. En proyectos reales, errores en estos cálculos pueden generar sobrecostos, retrasos y complicaciones técnicas que afectan tanto a los contratistas como a los usuarios finales de la infraestructura. Además, la correcta programación de actividades con un orden de precedencia lógico, apoyado en datos confiables y herramientas como MS Project, permite optimizar los recursos y garantizar una ejecución eficiente. Tal como se plantea más adelante, estos fueron los objetivos principales que se abordaron en la realización de este proyecto.


\section{Descripción del Proyecto}

El proyecto consite en la ampliación del Aeropuerto construyendo 
la extensión de la calle Diego Barros Ortiz y su respectivo paso 
bajo nivel, el cual cruza inferiormente a la pista de rodaje Mike.
Las Partidas incluidas en el proyecto son las siguientes:

\begin{itemize}
    \item \textbf{VIALIDAD}
    \begin{itemize}
        \item \textbf{Preparación área de trabajo}
        \begin{itemize}
            \item Remoción de Pavimento de Hormigón ($m^3$)
            \item Remoción de Pavimento Asfáltico ($m^3$)
        \end{itemize}
        \item \textbf{Movimiento de Tierra}
        \begin{itemize}
            \item Excavación en Terreno de Cualquier Naturaleza (TCN) ($m^3$) 
            \item Formación y COnstrucción de Terraplenes ($m^3$)
        \end{itemize}
        \item \textbf{Capas Granulares}
        \begin{itemize}
            \item Sub Base Granular CBR 40\% ($m^3$)
            \item Base Granular CBR 80\% ($m^3$)
            \item Carpeta Granular de Rodadura CBR 60\% ($m^3$)
        \end{itemize}
        \item \textbf{Pavimentos}
        \begin{itemize}
            \item Imprimación ($m^2$)
            \item Riego de Liga ($m^2$)
            \item Concreto Asfáltico, Capa de Rodadura ($m^3$)
            \item Concreto Asfáltico, Capa Intermedia ($m^3$)
            \item Aceras de Hormigón ($m^2$)
            \item Soleras Tipo A ($m$)
        \end{itemize}
        \item \textbf{Señalética y Demarcación}
        \begin{itemize}
            \item Señal Vertical Lateral Tipo 2 ($UN$)
            \item Señal Vertical Lateral Tipo 3 ($UN$)
            \item Líneas de Eje Segmentadas ($m$)
            \item Línea Continua ($m$)
            \item Líneas, Achurados, Símbolos y Leyendas ($m^2$)
        \end{itemize}
    \end{itemize}
    \item \textbf{ESTRUCTURAS}
    \begin{itemize}
        \item Excavación a Máquina en Puente y Estructuras ($m^3$)
        \item Relleno Estructural Permeable ($m^3$)
        \item Impermeabilización de Muros y Estribos ($m^2$)
        \item Hormigón G-5 ($m^3$)
        \item Hormigón G-30/SM ($m^3$)
        \item Acero para Armaduras A63-42H ($kg$)
        \item Moldajes para Infraestructura ($m^2$)
        \item Losa de Acceso ($m^3$)
        \item Barbacanas de desagüe ($UN$)
    \end{itemize}
\end{itemize}


\end{document}
